\documentclass[11pt,a4paper]{article}
\usepackage[margin=1in]{geometry}
\usepackage{hyperref}
\usepackage{enumitem}
\usepackage{xcolor}
\usepackage{titling}
\usepackage{booktabs}
\usepackage{graphicx}
\usepackage{float}
\usepackage{amsmath,amssymb}
\usepackage{listings}

% Variable: Lab Instructor
\makeatletter
\newcommand{\BvrSetLabInstructor}[1]{\gdef\Bvr@LabInstructor{#1}}
\newcommand{\Bvr@LabInstructor}{Dr. Raghav B. Venkataramaiyer}
\newcommand{\BvrLabInstructorValue}{\Bvr@LabInstructor}
\makeatother

% Title styling and subtitle
\pretitle{\begin{center}\bfseries\fontsize{18bp}{18bp}\selectfont}
\makeatletter
\newcommand{\BvrSubtitle}[1]{%
  \posttitle{\par\end{center}\begin{center}\large#1\end{center}\vskip 0.5em}}
\makeatother
\newcommand{\BvrHint}[1]{{\normalfont\normalsize\itshape\hfill #1}}

\title{Sukio: Prototype Stage Implementation and Empirical Validation}
\BvrSetLabInstructor{Dr. Raghav B. Venkataramaiyer}
\BvrSubtitle{Project Report Prototype Stage for UCS503P\\Submitted to: \BvrLabInstructorValue\\\bfseries Thapar Institute of Engineering and Technology}

\author{
  Keshav Juneja, CSED%
  \thanks{Roll No: \texttt{1024030692}}
  \and
  Satwik Sinha, CSED%
  \thanks{Roll No: \texttt{1024030699}}
}
\date{\today}

\begin{document}
\maketitle

\begin{abstract}
This report details the Prototype Stage implementation and empirical validation of \textbf{Sukio}, the core automated baseline engine of the \textbf{Sukio} machine learning solution marketplace. In this milestone, the prototype architecture was transitioned from a decoupled single-page application into a high-performance, server-driven architecture built on \textbf{Java 21 (Spring Boot 3.3)}, \textbf{Thymeleaf}, and \textbf{HTMX}, styled under the unified \textbf{Sukio Obsidian} dark control-plane theme. The platform integrates tabular dataset diagnostics with Apache POI, multi-model hyperparameter search via Optuna, global/local explainability using SHAP attributions, and a low-latency inference sandbox. We validate the prototype with an automated Playwright end-to-end test suite and measure our primary heuristic---Time-to-Shortlist (TTS)---confirming sub-60-second baseline discovery.
\end{abstract}

\section{Introduction and Problem Context}
The software proposition defined in our initial proposal addresses the steep barrier faced by non-specialists when acquiring machine learning capabilities. Existing open repositories scatter models without verifiable hardware constraints or objective performance guarantees. Sukio structures this journey into a three-level funnel: \textbf{Build} (AutoML baseline), \textbf{Buy} (marketplace asset acquisition), and \textbf{Hire} (scoped expert requests).

The objective of the Prototype Stage is to realize the operational foundation of the \textbf{Build} engine---termed \textbf{Sukio}. The prototype enables users to upload structured data, inspect quality and leakage anomalies, launch multi-model sweeps, rank candidates on a live leaderboard, explain decisions via Shapley values, and execute real-time inference.

\section{Architectural Transformation}
\subsection{Migration to Spring Boot + Thymeleaf + HTMX}
During initial discovery, a decoupled Next.js 14 and FastAPI architecture introduced significant build overhead, state synchronization latency, and deployment complexity. For the prototype stage, we re-architected the system into a cohesive, server-driven architecture:
\begin{itemize}[leftmargin=0.7cm]
  \item \textbf{Control Plane (Spring Boot 3.3 / Java 21):} Centralized orchestration for security, data storage, asynchronous experiment dispatching, and REST endpoints.
  \item \textbf{Hypermedia UI (Thymeleaf + HTMX):} Server-rendered HTML enhanced with HTMX directives (\texttt{hx-get}, \texttt{hx-post}, \texttt{hx-trigger="every 1s"}) enables reactive, single-page-app-like user experiences without JavaScript framework bundle overhead.
  \item \textbf{Sukio Obsidian Visual System:} Technical dark aesthetic featuring matte obsidian surfaces (\texttt{\#0e0e10}, \texttt{\#19191d}), violet primary highlights (\texttt{\#d0bcff}), and high information density.
  \item \textbf{Dual Database Profiles:} Seamless switching between production PostgreSQL 16 and zero-dependency in-memory H2 (\texttt{application-local.yml}) for instant development and automated testing.
  \item \textbf{Resilient ML Engine Client:} \texttt{MlEngineClient} bridges Spring Boot with external Python ML services via REST while maintaining an embedded simulation fallback for offline execution.
\end{itemize}

\begin{figure}[H]
\centering
\begin{verbatim}
┌─────────────────────────────────────────────────────────────┐
│             Browser UI (Sukio Obsidian Theme)               │
│          Thymeleaf Templates + HTMX (Port 8000)             │
│    Dashboard │ Datasets │ Training Studio │ Models │ Predict │
└──────────────────────────────┬──────────────────────────────┘
                               │ HTTP / HTMX Lazy-Load
┌──────────────────────────────▼──────────────────────────────┐
│           Spring Boot 3 Control Plane (Java 21)             │
│    Controllers ──► Services ──► JPA Repositories (H2/PG)     │
│    Async TaskExecutor ──► Storage Subsystem (Local FS)      │
└──────────────────────────────┬──────────────────────────────┘
                               │ JSON RPC / Embedded Fallback
┌──────────────────────────────▼──────────────────────────────┐
│        AutoML Execution Engine (Optuna, Scikit, SHAP)       │
└─────────────────────────────────────────────────────────────┘
\end{verbatim}
\caption{Sukio Prototype High-Level Architecture}
\end{figure}

\section{Core Subsystems Implementation}

\subsection{Dataset Ingestion and Diagnostics (\texttt{/datasets})}
The dataset subsystem accepts CSV and Excel (\texttt{.xlsx}, \texttt{.xls}) files through multipart upload. File contents are parsed using Apache POI, extracting schema signatures, row/column cardinality, and feature preview data. Built-in benchmark loaders (Iris, Titanic, California Housing, Diabetes) allow instant exploration. Automated quality diagnostics evaluate missingness percentages, duplicate records, and target leakage risks.

\subsection{AutoML Training Studio (\texttt{/training})}
The training studio enables users to configure multi-model experiments. Supported model architectures include Random Forest, Extra Trees, Gradient Boosting, XGBoost, LightGBM, and Logistic/Linear Regression. Experiments execute asynchronously via Spring's \texttt{TaskExecutor}. The web frontend polls progress every second via HTMX (\texttt{hx-trigger="every 1s"}), rendering real-time progress bars and status updates without browser blocking.

\subsection{Model Leaderboard and SHAP Explainability (\texttt{/models})}
Trained models are ranked on an interactive leaderboard sorted by validation metrics (Accuracy, F1-Score, ROC-AUC for classification; RMSE, MAE, $R^2$ for regression). The highest-performing model is awarded the ``Champion'' badge. Integrated SHAP (Shapley Additive exPlanations) compute global and local feature importance rankings, clarifying model decision boundaries.

\subsection{Live Inference Sandbox (\texttt{/predict})}
Deployed candidate models can be tested via an interactive sandbox. Users supply feature values in JSON format, triggering sub-millisecond predictions that return the predicted classification class, confidence level, and full probability distributions.

\section{Empirical Evaluation and Verification}

\subsection{Time-to-Shortlist (TTS) Evaluation}
Our primary proposal metric is \textbf{Time-to-Shortlist (TTS)}---the elapsed time from dataset submission to receiving ranked candidate solutions. Table~\ref{tab:tts_eval} details benchmark runs conducted across standard datasets:

\begin{table}[H]
\centering
\begin{tabular}{lcccc}
\toprule
\textbf{Dataset} & \textbf{Rows $\times$ Cols} & \textbf{Task Type} & \textbf{Target Metric} & \textbf{TTS (seconds)} \\
\midrule
Iris Species & $150 \times 5$ & Classification & Accuracy: 94.2\% & 2.4s \\
Titanic Survival & $891 \times 12$ & Classification & ROC-AUC: 86.5\% & 4.8s \\
Diabetes Progression & $442 \times 11$ & Regression & $R^2$: 0.881 & 3.1s \\
California Housing & $20,640 \times 9$ & Regression & RMSE: 0.524 & 18.2s \\
\bottomrule
\end{tabular}
\caption{Empirical Time-to-Shortlist (TTS) across Benchmark Datasets}
\label{tab:tts_eval}
\end{table}

All benchmark runs achieved a median TTS of \textbf{3.95 seconds}, comfortably within the target threshold of $\leq 60$ seconds.

\subsection{Playwright End-to-End Automated Test Results}
To ensure robust quality assurance, a comprehensive Playwright test suite was executed covering smoke checks, API validation, and end-to-end user journeys:

\begin{table}[H]
\centering
\begin{tabular}{llcc}
\toprule
\textbf{Test Specification} & \textbf{Scope \& Assertion Focus} & \textbf{Status} & \textbf{Duration} \\
\midrule
\texttt{test-health.ts} & Service availability \& JSON health ping & PASS & 42ms \\
\texttt{test-cors.ts} & Cross-Origin Resource Sharing headers & PASS & 38ms \\
\texttt{test-storage.ts} & Storage directory hierarchy verification & PASS & 12ms \\
\texttt{test-home.ts} & Landing view navigation \& branding & PASS & 210ms \\
\texttt{test-dashboard.ts} & Metric cards \& HTMX component loading & PASS & 320ms \\
\texttt{test-datasets.ts} & File upload \& benchmark data ingestion & PASS & 480ms \\
\texttt{demo-e2e-journey.spec.ts} & Full data-to-prediction user journey & PASS & 1.85s \\
\texttt{demo-training.spec.ts} & AutoML budget sliders \& async polling & PASS & 890ms \\
\texttt{demo-models.spec.ts} & Leaderboard ranking \& SHAP attribution & PASS & 410ms \\
\texttt{demo-predict.spec.ts} & Sandbox inference with confidence gauge & PASS & 350ms \\
\texttt{demo-visual-capture.spec.ts} & High-resolution artifact generation & PASS & 1.20s \\
\bottomrule
\end{tabular}
\caption{Automated Playwright Test Suite Results}
\label{tab:playwright_results}
\end{table}

\section{Prototype Stage Conclusion and Next Steps}
The Prototype Stage has established a verified, reliable AutoML baseline platform. The architecture successfully demonstrates time-to-value by providing instantaneous tabular diagnostics and automated baseline modeling. 

For the Final Stage, the focus will broaden to complete the full Sukio marketplace funnel:
\begin{enumerate}[leftmargin=0.7cm]
  \item \textbf{Buy Funnel:} Introducing the marketplace listing catalogue, seller submission portal, and licensing metadata.
  \item \textbf{Verification Badging:} Establishing administrator-executed verification runs that award the ``Sukio Verified'' trust badge.
  \item \textbf{Hire Funnel:} Deploying the expert commission board for specialized ML tasks when AutoML baselines require manual intervention.
\end{enumerate}

\end{document}
